\clearpage % Rozdziały zaczynamy od nowej strony.
\section{Wstęp}

Rozwój elektroniki pojazdowej sprawił, że współczesny samochód dysponuje dziesiątkami
czujników i elektronicznych jednostek sterujących (ECU), które w czasie rzeczywistym
rejestrują parametry pracy silnika, układu hamulcowego, przekładni oraz wielu innych
podzespołów. Dane te, przesyłane najczęściej magistralą CAN~\cite{bosch1991}, od dawna
wykorzystywane są przez producentów pojazdów do diagnostyki i sterowania, jednak
w ostatnich latach coraz częściej stają się też podstawą analizy samego kierowcy ---
sposobu, w jaki prowadzi on pojazd. Styl jazdy ma bowiem bezpośredni, mierzalny wpływ na
koszty eksploatacji, emisję spalin oraz bezpieczeństwo ruchu drogowego. Badania
przeprowadzone przez Oak Ridge National Laboratory na zlecenie amerykańskiego
Departamentu Energii wykazały, że agresywny styl jazdy --- gwałtowne przyspieszanie
i hamowanie --- może zwiększyć zużycie paliwa nawet o 10--40\% w ruchu miejskim oraz
o 15--30\% przy prędkościach charakterystycznych dla tras szybkiego ruchu, w porównaniu
do jazdy spokojnej i płynnej~\cite{thomas2017}. Skala tego zjawiska sprawia, że
informowanie kierowcy o jego stylu jazdy oraz motywowanie go do jazdy ekonomicznej
i bezpiecznej stało się osobnym, aktywnie rozwijanym kierunkiem zarówno w przemyśle
motoryzacyjnym, jak i w badaniach naukowych.

Równolegle do rozwoju samej elektroniki pojazdowej dynamicznie rozwija się rynek
systemów telematycznych --- rozwiązań łączących zbieranie danych z pojazdu, ich
transmisję oraz analizę, najczęściej z wykorzystaniem chmury obliczeniowej.
Zastosowania takich systemów obejmują zarządzanie flotami pojazdów, ubezpieczenia
komunikacyjne taryfikowane na podstawie rzeczywistego stylu jazdy (UBI), w których
wysokość składki zależy bezpośrednio od zarejestrowanych podczas jazdy parametrów,
a nie wyłącznie od statystycznych cech kierowcy takich jak wiek, staż czy dotychczasowa
historia szkód, a także aplikacje konsumenckie, których celem jest poprawa nawyków
kierowcy. Wspólnym
mianownikiem większości dostępnych na rynku rozwiązań jest jednak zależność od
łączności sieciowej oraz przetwarzania danych poza pojazdem --- na serwerach dostawcy
usługi. Ocena stylu jazdy jest w takich systemach zwykle wynikiem analizy wykonywanej
po zakończeniu przejazdu lub w chmurze w czasie zbliżonym do rzeczywistego, rzadziej zaś
efektem wnioskowania realizowanego bezpośrednio na pokładzie pojazdu, bez udziału
zewnętrznej infrastruktury. Taka architektura, choć wygodna z punktu widzenia dostawcy
usługi (możliwość centralnego gromadzenia danych, aktualizacji algorytmów, rozliczeń
abonamentowych), wiąże funkcjonowanie urządzenia pokładowego z ciągłą dostępnością
łączności komórkowej oraz z istnieniem samej usługi w czasie --- co, jak pokazano
w dalszej części tego rozdziału, w praktyce bywa istotnym ograniczeniem.

Innym istotnym nurtem badawczym jest wykorzystanie metod uczenia maszynowego do
automatycznej klasyfikacji stylu jazdy na podstawie danych pojazdowych lub danych
z czujników inercyjnych. Podejścia takie pozwalają zastąpić proste, progowe reguły
(np. wykrywanie przekroczenia zadanej wartości przyspieszenia) modelami zdolnymi do
jednoczesnego uwzględnienia wielu, często wzajemnie skorelowanych cech jazdy, co
przekłada się na trafniejszą i bardziej odporną na przypadkowe zdarzenia ocenę
zachowania kierowcy~\cite{meiring2015}. Równocześnie coraz więcej uwagi poświęca się nie
tylko samej ocenie stylu jazdy, ale też sposobowi przekazania jej kierowcy. Badania
z zakresu ekonomii behawioralnej i projektowania gier wskazują, że mechanizmy
motywacyjne, nagradzające pożądane zachowania, są skuteczniejszą formą oddziaływania na
nawyki niż komunikaty ostrzegawcze i informacje o błędach~\cite{thaler2008,deterding2011}.
Połączenie tych dwóch obserwacji --- że jakość danych wejściowych i metoda ich analizy
determinują trafność oceny, a forma przekazania wyniku determinuje jej rzeczywisty
wpływ na zachowanie kierowcy --- stanowi punkt wyjścia dla rozwiązania opisanego
w niniejszej pracy.

Opisane w niniejszej pracy rozwiązanie łączy trzy elementy omówione powyżej:
wykorzystanie rzeczywistych danych pojazdowych z magistrali CAN (a nie danych
z czujników smartfona), klasyfikację stylu jazdy metodą uczenia maszynowego (a nie
prostych reguł progowych) oraz przekazanie wyniku kierowcy w formie zgamifikowanej,
motywującej informacji zwrotnej (a nie liczbowego wskaźnika lub komunikatu
ostrzegawczego) --- realizowane w całości lokalnie, na pokładzie pojazdu, bez
konieczności stałej łączności sieciowej ani przesyłania danych na zewnątrz. W dalszej
części niniejszego rozdziału przedstawiono przegląd wybranych rozwiązań dostępnych
obecnie na rynku, naukowe uzasadnienie przyjętych
w pracy metod, a następnie omówiono
proponowane w pracy rozwiązanie wraz z jego przewagami względem opisanych alternatyw.

\subsection{Przegląd istniejących rozwiązań}
\label{sec:przeglad}

Systemy analizujące styl jazdy kierowcy można podzielić na kilka kategorii, różniących
się przede wszystkim źródłem danych (czujniki smartfona, złącze OBD, bezpośredni dostęp
do magistrali CAN), miejscem przetwarzania (chmura, aplikacja mobilna, urządzenie
pokładowe) oraz sposobem prezentacji wyniku kierowcy (liczbowy wskaźnik, ranking,
informacja zwrotna w czasie rzeczywistym). W niniejszym podrozdziale omówiono trzy
reprezentatywne przykłady z tych kategorii: komercyjną platformę telematyczną opartą
na czujnikach smartfona, historyczny, lecz wciąż często przywoływany w literaturze
system wykorzystujący dane z sieci pokładowej pojazdu, a także rozwiązania telematyki
flotowej i ubezpieczeniowej.

\subsubsection{Platformy telematyczne oparte na czujnikach smartfona}

Największą grupę wdrożeń komercyjnych stanowią dziś systemy wykorzystujące wyłącznie
czujniki wbudowane w smartfon kierowcy --- akcelerometr, żyroskop oraz odbiornik GPS ---
bez potrzeby instalowania jakiegokolwiek dodatkowego sprzętu w pojeździe. Przykładem
takiego podejścia jest platforma DriveWell firmy Cambridge Mobile Telematics,
wykorzystywana przez wielu ubezpieczycieli do wdrażania ubezpieczeń UBI. Aplikacja
instalowana na telefonie kierowcy zbiera dane z czujników ruchu i lokalizacji, na ich
podstawie wykrywa m.in. gwałtowne hamowanie, przekraczanie prędkości oraz rozproszenie
uwagi kierowcy wynikające z korzystania z telefonu w trakcie jazdy, a następnie
przypisuje przejazdowi ocenę w skali od 1 do 100~\cite{cmtelematics}. Wyniki są
przekazywane kierowcy w formie punktowej oraz motywacyjnej (system nagród za dobre
przejazdy), a producent deklaruje, że regularne korzystanie z aplikacji ogranicza
rozproszenie uwagi kierowcy o ok. 35\%, a przypadki przekraczania prędkości i
gwałtownego hamowania o ok. 20\%~\cite{cmtelematics}. Platforma umożliwia też
opcjonalne podłączenie dedykowanych urządzeń (tzw. Tags) w celu podniesienia
dokładności pomiaru, co producent wprost wskazuje jako sposób na ograniczenie
niedokładności wynikającej z samych tylko czujników telefonu~\cite{cmtelematics}.
Rozwiązanie to ewoluowało w kierunku łączenia (fuzji) danych z wielu niezależnych
źródeł jednocześnie --- oprócz smartfona i wspomnianych urządzeń Tags, platforma
DriveWell obsługuje również dane z rejestratorów jazdy (dashcam) oraz z modułów
łączności wbudowanych fabrycznie w nowsze pojazdy, a mechanizm automatycznej
identyfikacji kierowcy i pasażerów pozwala przypisać zarejestrowany przejazd
właściwej osobie bez ręcznej konfiguracji~\cite{cmtelematics}. Takie podejście dobrze
ilustruje ogólny kierunek rozwoju komercyjnych platform telematycznych: zamiast
pojedynczego źródła danych, systemy te dążą do łączenia wielu niezależnych, z osobna
niedoskonałych źródeł w celu skompensowania ograniczeń każdego z nich.

Zaletą tego podejścia jest brak konieczności ingerencji w instalację elektryczną
pojazdu oraz niski koszt wdrożenia --- jedynym wymaganym elementem jest smartfon,
którym dysponuje niemal każdy kierowca. Istotnym ograniczeniem jest jednak jakość
i wiarygodność danych źródłowych: pomiary przyspieszenia wykonywane czujnikami
smartfona zależą od sposobu jego zamocowania w pojeździe, mogą być zakłócone przez ruch
samego telefonu (np. swobodnie leżącego na siedzeniu) i nie odzwierciedlają bezpośrednio
parametrów pracy pojazdu, takich jak rzeczywista prędkość obrotowa silnika, położenie
pedału przyspieszenia czy ciśnienie w układzie hamulcowym, dostępnych natywnie na
magistrali CAN. Ocena stylu jazdy odbywa się ponadto po stronie serwera dostawcy usługi,
co oznacza zależność od łączności sieciowej oraz przesyłanie danych o lokalizacji
i sposobie jazdy poza pojazd.

\subsubsection{System Fiat eco:Drive}

Interesującym, choć już historycznym przykładem systemu wykorzystującego bezpośrednio
dane z sieci pokładowej pojazdu jest aplikacja eco:Drive, zaprezentowana przez Fiata
na salonie samochodowym w Paryżu w 2008 roku~\cite{fiatecodrive}. W odróżnieniu od
rozwiązań opartych na smartfonie, eco:Drive korzystała z danych rejestrowanych
bezpośrednio w sieci pokładowej pojazdu, obejmujących m.in. użyty bieg, obroty silnika
oraz sposób korzystania z pedałów przyspieszenia i hamulca, zapisywanych na karcie
pamięci USB i wgrywanych po zakończeniu jazdy do aplikacji na komputerze
osobistym~\cite{fiatecodrive}. Aplikacja wskazywała kierowcy konkretne fragmenty
przejazdu, w których styl jazdy odbiegał od ekonomicznego (np. zbyt intensywne
przyspieszanie lub jazda na zbyt niskim biegu), a w rozszerzonej wersji eco:Drive
Social, zaprezentowanej podczas salonu samochodowego we Frankfurcie, wprowadzono
elementy gamifikacji oparte na dynamice społecznościowej --- odznaki potwierdzające
dobre nawyki za kierownicą oraz ranking społecznościowy porównujący kierowców między
sobą~\cite{fiatecodrivesocial}.
Istotną cechą eco:Drive, odróżniającą ją od generycznych porad dotyczących
oszczędnej jazdy, była szczegółowość informacji zwrotnej: aplikacja nie ograniczała się
do ogólnego zalecenia w rodzaju \emph{jeźdź spokojniej}, lecz precyzyjnie wskazywała,
w którym konkretnie przejeździe i w którym jego fragmencie kierowca przyspieszał zbyt
gwałtownie lub jechał na niewłaściwym biegu~\cite{fiatecodrive}. Aplikację, dostępną na
wielu platformach dzięki wykorzystaniu środowiska Adobe Air, opracowała na zlecenie
Fiata brytyjska agencja konsultingowa AKQA~\cite{fiatecodrive}.
Zebrane przez Fiata dane z rzeczywistej eksploatacji wykazały średnio 6-procentową
redukcję zużycia paliwa wśród użytkowników aplikacji, sięgającą 16\% w grupie 10\%
kierowców, którzy najbardziej poprawili swój styl jazdy~\cite{fiatecodrive}.

System ten pozostaje istotnym punktem odniesienia dla niniejszej pracy z dwóch powodów:
wykorzystuje dane pojazdowe analogiczne do wykorzystywanych w opisywanym dalej
rozwiązaniu, a jednocześnie stosuje elementy gamifikacji zamiast prostych komunikatów
ostrzegawczych. Jego zasadniczym ograniczeniem jest jednak brak informacji zwrotnej
w czasie rzeczywistym --- analiza i wizualizacja wyników odbywały się dopiero po
zakończeniu jazdy, na osobnym urządzeniu (komputerze osobistym), co znacząco ogranicza
możliwość bieżącej korekty zachowania za kierownicą w trakcie samej jazdy.

\subsubsection{Telematyka flotowa i ubezpieczeniowa}

Odrębną kategorię stanowią systemy telematyki flotowej, projektowane z myślą
o zarządcach dużych flot pojazdów, a nie o indywidualnym kierowcy. Przykładem jest
platforma firmy Geotab, w której urządzenie podłączane do złącza OBD rejestruje
i przesyła do chmury nawet ponad 250 punktów danych na sekundę, obejmujących
m.in. przekroczenia prędkości, gwałtowne hamowanie, gwałtowne przyspieszanie oraz
gwałtowne pokonywanie zakrętów~\cite{geotab}. Zdarzenia te wykrywane są na podstawie
progów wartości przyspieszenia skalibrowanych względem klasy pojazdu, a wynikowa ocena
bezpieczeństwa jazdy prezentowana jest zarządcy floty w panelu przeglądarkowym, wraz
z możliwością konfiguracji alertów w czasie rzeczywistym~\cite{geotab}. Tak zebrane dane
wykorzystywane są przede wszystkim do coachingu kierowców przez zarządcę floty oraz do
budowania rankingów porównujących poszczególnych kierowców lub całe floty względem
grupy odniesienia --- producent przywołuje przykład floty, która dzięki takiemu
podejściu podniosła swoją średnią ocenę bezpieczeństwa jazdy do 88 punktów na 100,
uzyskując wynik o 6\% lepszy niż średnia najlepszych flot z tej samej grupy
porównawczej~\cite{geotab}. W odróżnieniu od podejść opisanych wcześniej, ocena stylu
jazdy opiera się tu jednak na prostych regułach progowych, a nie na modelu uczenia
maszynowego, zaś odbiorcą informacji zwrotnej jest przede wszystkim zarządca floty,
nie sam kierowca prowadzący pojazd w danej chwili.

Zbliżoną architekturę, adresowaną jednak do klienta indywidualnego, prezentowało
urządzenie Automatic --- podłączany do złącza OBD moduł komunikujący się bezprzewodowo
ze smartfonem, który udostępniał m.in. powiadomienia o wypadku, historię przejazdów,
diagnostykę pojazdu oraz cotygodniową, zbiorczą ocenę stylu jazdy, wyliczaną na
podstawie zdarzeń zarejestrowanych i przetworzonych przez producenta usługi w chmurze.
Wokół urządzenia zbudowano ponadto ekosystem integracji z aplikacjami zewnętrznymi,
umożliwiający m.in. automatyczne rejestrowanie przejazdów służbowych na potrzeby
rozliczeń podatkowych czy udostępnianie danych o lokalizacji pojazdu innym usługom ---
co w praktyce oznaczało, że wartość urządzenia dla użytkownika wykraczała poza samą
ocenę stylu jazdy i obejmowała całą warstwę usług zbudowanych wokół danych z pojazdu.
Usługa ta, mimo dużej popularności, została jednak całkowicie wyłączona w 2020~roku,
a wraz z nią przestały działać wszystkie należące do niej urządzenia --- decyzję tę
producent uzasadnił brakiem rentowności usługi w obliczu spadku sprzedaży i wynajmu
pojazdów w okresie pandemii~\cite{automaticlabs}. Przypadek ten dobrze ilustruje istotną
wadę modelu, w którym funkcjonalność urządzenia pokładowego jest trwale uzależniona od
dostępności usługi chmurowej dostawcy: zaprzestanie działania serwera dostawcy oznacza
całkowitą utratę funkcjonalności sprzętu, niezależnie od jego sprawności technicznej.

\subsubsection{Analiza wad i zalet istniejących rozwiązań}

Zestawienie omówionych rozwiązań pod względem źródła danych, miejsca ich przetwarzania
oraz formy informacji zwrotnej przekazywanej kierowcy przedstawiono w
tabeli~\ref{tab:analiza-rozwiazan}. Wynika z niego, że żadne z analizowanych rozwiązań
nie łączy jednocześnie danych z magistrali CAN, przetwarzania lokalnego (bez zależności
od łączności sieciowej) oraz gamifikowanej informacji zwrotnej dostępnej w czasie
rzeczywistym. Rozwiązanie oparte na czujnikach smartfona (DriveWell) cierpi na
niedokładność pomiaru wynikającą z pośredniego charakteru danych. Rozwiązania
wykorzystujące rzeczywiste dane pojazdowe (eco:Drive, Geotab) nie oferują natomiast
informacji zwrotnej w czasie rzeczywistym bezpośrednio w pojeździe --- analiza odbywa
się po zakończeniu jazdy (eco:Drive) albo trafia wyłącznie do zarządcy floty, z pominięciem
samego kierowcy (Geotab). Wszystkie trzy omówione rozwiązania zależne są ponadto od
chmury obliczeniowej dostawcy usługi, co wiąże trwałość działania urządzenia
z istnieniem tej usługi w czasie --- co w przypadku Automatic doprowadziło do
całkowitej utraty funkcjonalności sprzętu użytkowników po wycofaniu usługi przez
producenta.

Drugą, niezależną od źródła danych osią podziału jest moment, w którym kierowca
otrzymuje informację zwrotną o swoim stylu jazdy. We wszystkich trzech omówionych
rozwiązaniach informacja ta dociera do kierowcy dopiero po zakończeniu przejazdu ---
w formie zbiorczej oceny liczbowej (DriveWell, Automatic), raportu analizowanego po
podłączeniu urządzenia do komputera (eco:Drive) lub w ogóle nie dociera bezpośrednio
do kierowcy, lecz do osoby trzeciej zarządzającej flotą (Geotab). Taki retrospektywny
charakter informacji zwrotnej ogranicza możliwość natychmiastowej korekty zachowania
w chwili, w której do niekorzystnego zdarzenia faktycznie dochodzi --- kierowca
dowiaduje się o popełnionym błędzie dopiero po fakcie, gdy nie ma już możliwości
bezpośredniego powiązania oceny z konkretną decyzją podjętą za kierownicą.

\begin{table}[!h]
    \centering
    \caption{Zestawienie porównawcze przeglądniętych rozwiązań analizy stylu jazdy.}
    \label{tab:analiza-rozwiazan}
    \begin{tabular}{| p{2.6cm} | p{2.6cm} | p{2.6cm} | p{2.5cm} | p{3.0cm} |} \hline
        \textbf{Rozwiązanie} & \textbf{Źródło danych} & \textbf{Metoda oceny} & \textbf{Miejsce przetwarzania} & \textbf{Informacja zwrotna} \\ \hline\hline
        DriveWell (CMT)~\cite{cmtelematics} & czujniki smartfona (GPS, akcelerometr) & reguły + modele w~chmurze & chmura dostawcy usługi & wynik liczbowy (1--100), po zakończeniu jazdy \\ \hline
        Fiat eco:Drive~\cite{fiatecodrive,fiatecodrivesocial} & magistrala pojazdu (bieg, obroty, pedały) & reguły progowe & aplikacja PC (po jeździe) & gamifikacja (odznaki, ranking), po zakończeniu jazdy \\ \hline
        Geotab~\cite{geotab} & OBD (przyspieszenie, prędkość) & reguły progowe & chmura dostawcy usługi & panel dla zarządcy floty, nie dla kierowcy \\ \hline
        Automatic~\cite{automaticlabs} & OBD + smartfon & reguły w~chmurze & chmura dostawcy usługi (usługa wycofana) & ocena tygodniowa w~aplikacji \\ \hline
        \textbf{Proponowane rozwiązanie} & \textbf{magistrala CAN} & \textbf{uczenie maszynowe} & \textbf{lokalnie, na~pokładzie pojazdu} & \textbf{gamifikacja w~czasie rzeczywistym} \\ \hline
    \end{tabular}
\end{table}

\subsection{Naukowe uzasadnienie przyjętych metod}
\label{sec:uzasadnienie-metod}

Niezależnie od rozwiązań komercyjnych omówionych w podrozdziale~\ref{sec:przeglad},
zastosowanie w niniejszej pracy dwóch rozwiązań metodycznych --- klasyfikacji stylu
jazdy za pomocą algorytmów uczenia maszynowego oraz gamifikacji jako formy prezentacji
wyniku tej klasyfikacji kierowcy --- znajduje uzasadnienie w wynikach badań
naukowych publikowanych w literaturze przedmiotu. W dalszej części podrozdziału
scharakteryzowano stan wiedzy w obu tych obszarach, przywołując konkretne wyniki
badawcze potwierdzające skuteczność każdego z tych podejść.

\subsubsection{Klasyfikacja stylu jazdy metodami uczenia maszynowego}

Skuteczność klasyfikacji stylu jazdy metodami uczenia maszynowego potwierdza bogata
literatura naukowa. Meiring i Myburgh w obszernym przeglądzie systematyzują stosowane
w tym obszarze metody sztucznej inteligencji, wskazując na rosnącą przewagę podejść
opartych na uczeniu maszynowym nad prostymi regułami progowymi, wynikającą z ich
zdolności do jednoczesnego uwzględnienia wielu, często wzajemnie skorelowanych cech
jazdy~\cite{meiring2015}. W ramach przeglądu autorzy wskazują w szczególności na
obiecujący potencjał metod opartych na logice rozmytej, ukrytych modelach Markowa
oraz maszynach wektorów nośnych, zastrzegając jednocześnie, że praktyczne
zastosowanie tych metod do jednoznacznej identyfikacji indywidualnego kierowcy wymaga
dalszego ograniczenia złożoności obliczeniowej modeli~\cite{meiring2015}. Bejani
i Ghatee zaproponowali system oceny stylu jazdy wykorzystujący uczenie zespołowe (ang.
\emph{ensemble learning}) na danych z czujników smartfona, uwzględniający dodatkowo
kontekst jazdy (np.\ rodzaj drogi), co pozwoliło ograniczyć liczbę błędnych klasyfikacji
wynikających z bezpośredniego porównywania jazdy miejskiej i autostradowej według tych
samych kryteriów~\cite{bejani2018}. W zaproponowanym przez nich podejściu, opartym na
logice neuronowo-rozmytej, styl jazdy klasyfikowany jest na podstawie podobieństwa
zarejestrowanego przebiegu do wzorcowych profili rozmytych, wyznaczonych statystycznie
na podstawie danych z czujników smartfona~\cite{bejani2018}. Khan
i inni wykorzystali z kolei dane pochodzące bezpośrednio z magistrali CAN pojazdu do
identyfikacji kierowcy metodami uczenia maszynowego, potwierdzając, że parametry
rejestrowane natywnie przez sieć pokładową pojazdu (m.in.\ prędkość, obroty silnika,
położenie przepustnicy) niosą wystarczająco dużo informacji do skutecznej klasyfikacji
zachowań za kierownicą, bez konieczności sięgania po dodatkowe czujniki
inercyjne~\cite{khan2022}.

Przywołane prace potwierdzają dwa istotne dla niniejszej pracy założenia: po pierwsze,
że modele uczenia maszynowego skuteczniej niż proste reguły progowe radzą sobie
z klasyfikacją stylu jazdy dzięki zdolności do jednoczesnego uwzględnienia wielu
skorelowanych cech, a po drugie, że dane rejestrowane natywnie na magistrali CAN
pojazdu stanowią samodzielnie wystarczające i wiarygodne źródło do takiej klasyfikacji.
Na tych właśnie założeniach oparto wybór metody klasyfikacji zastosowanej
w opracowywanym systemie, opisanej szczegółowo w rozdziale~\ref{sec:trening}.

\subsubsection{Gamifikacja informacji zwrotnej o stylu jazdy}

Drugim filarem metodycznym niniejszej pracy jest przekazanie wyniku klasyfikacji
kierowcy w formie zgamifikowanej. Teoretyczne podstawy takiego podejścia wywodzą się
z koncepcji nudgingu, sformułowanej przez Thalera i Sunsteina w ramach tzw.
paternalizmu libertariańskiego. Nudging definiowany jest jako dowolny element
architektury wyboru, który w przewidywalny sposób modyfikuje zachowanie ludzi, nie
zakazując przy tym żadnej z dostępnych opcji ani nie zmieniając w istotny sposób
bodźców ekonomicznych związanych z podejmowaną decyzją~\cite{thaler2008}. Klasycznym
przykładem takiego mechanizmu jest zmiana domyślnej opcji wyboru --- np.\ ustawienie
domyślnego uczestnictwa w programie emerytalnym jako \emph{aktywne} zamiast
\emph{nieaktywne} istotnie zwiększa liczbę uczestników programu, mimo że możliwość
rezygnacji pozostaje w pełni dostępna~\cite{thaler2008}. W odróżnieniu
od zakazów, kar czy komunikatów ostrzegawczych, nudge oddziałuje na decyzję poprzez
zmianę sposobu prezentacji informacji lub kontekstu, w którym decyzja jest
podejmowana --- np.\ poprzez natychmiastową, pozytywną informację zwrotną o pożądanym
zachowaniu zamiast sygnalizowania błędów. Pokrewnym pojęciem, wykorzystanym
równolegle w niniejszej pracy, jest gamifikacja, definiowana jako zastosowanie
elementów typowo kojarzonych z grami (punktacji, poziomów zaawansowania, odznak,
rankingów) w kontekście niezwiązanym z rozrywką, w celu zwiększenia zaangażowania
użytkownika~\cite{deterding2011}. Oba pojęcia, choć wywodzące się z różnych
dziedzin --- ekonomii behawioralnej oraz projektowania gier --- opierają się na
wspólnym założeniu, że motywowanie poprzez nagradzanie pożądanych zachowań jest
skuteczniejszą formą oddziaływania na nawyki niż komunikaty ostrzegawcze i informacje
o błędach~\cite{thaler2008,deterding2011}. Założenie to znajduje potwierdzenie również
w badaniach dotyczących bezpośrednio stylu jazdy. Vitiello i inni przeprowadzili
pilotażowe badanie rzeczywistego wdrożenia aplikacji mobilnej łączącej analizę stylu
jazdy w czasie rzeczywistym z elementami gamifikacji, potwierdzając pozytywny wpływ
takiego podejścia na trwałość zmiany nawyków kierowców w porównaniu z przekazywaniem
samej tylko oceny liczbowej~\cite{vitiello2023}. Podobne podejście zaprezentowali
Avolicino i inni w systemie EcoGO, w którym ranking, punktacja oraz poziomy
zaawansowania towarzyszą informacji zwrotnej przekazywanej kierowcy w czasie
rzeczywistym na podstawie prędkości oraz intensywności przyspieszania
i hamowania~\cite{avolicino2022}. Obie prace potwierdzają empirycznie tezę formułowaną
na gruncie ekonomii behawioralnej i projektowania gier, a jednocześnie pokazują, że
sama poprawa stylu jazdy przekłada się wymiernie na ograniczenie zużycia paliwa
i emisji spalin~\cite{barkenbus2010}.

Przywołane badania, podobnie jak wcześniej opisany przykład systemu Fiat
eco:Drive~(podrozdział~\ref{sec:przeglad}), potwierdzają zasadność zastosowania
gamifikacji w interfejsie opracowywanego systemu. W obu przywołanych pracach
gamifikowana informacja zwrotna realizowana jest w formie osobnej aplikacji mobilnej,
korzystającej z danych czujników smartfona --- w niniejszej pracy to samo podejście
motywacyjne zastosowano natomiast w oparciu o dane z magistrali CAN pojazdu oraz
zintegrowany z nim, a nie zewnętrzny, wyświetlacz, co szczegółowo opisano
w podrozdziale~\ref{sec:ui}.

\subsection{Proponowane rozwiązanie}
\label{sec:proponowane}

Na podstawie przeprowadzonego przeglądu sformułowano założenia rozwiązania opisanego
szczegółowo w rozdziale~\ref{sec:architektura} niniejszej pracy. Zgodnie z celem pracy
przedstawionym we wcześniejszej części dokumentu, proponowany system zbiera dane
bezpośrednio z magistrali CAN pojazdu, klasyfikuje styl jazdy kierowcy metodą uczenia
maszynowego oraz przekazuje wynik tej klasyfikacji w formie zgamifikowanej,
motywującej informacji zwrotnej, wyświetlanej w czasie rzeczywistym na zintegrowanym
z urządzeniem ekranie dotykowym. Całość przetwarzania --- od akwizycji danych, przez
klasyfikację modelem uczenia maszynowego, po generowanie informacji zwrotnej ---
odbywa się lokalnie, na pojedynczej platformie sprzętowej zamontowanej w pojeździe, bez
udziału zewnętrznych serwerów, aplikacji mobilnej ani stałej łączności sieciowej.

\subsubsection{Unikalne cechy proponowanego rozwiązania}

Proponowane rozwiązanie wyróżnia na tle opisanych wcześniej systemów kilka cech.
Po pierwsze, źródłem danych wejściowych jest bezpośrednio magistrala CAN pojazdu, a nie
czujniki inercyjne smartfona ani ograniczony podzbiór parametrów odczytywanych ze
złącza OBD. Pozwala to na wykorzystanie tych samych, precyzyjnych sygnałów, z których
korzystają jednostki sterujące samego pojazdu (rzeczywista prędkość, obroty silnika,
położenie pedału przyspieszenia, stan hamulca), bez niepewności pomiarowej wynikającej
z pośredniego charakteru pomiaru czujnikami niezwiązanymi na stałe z pojazdem. Po
drugie, klasyfikacja stylu jazdy wykonywana jest lokalnym modelem uczenia maszynowego,
wytrenowanym w sposób opisany w rozdziale~\ref{sec:trening} i wyeksportowanym do postaci
statycznego kodu uruchamianego bezpośrednio na mikrokontrolerze platformy docelowej ---
w odróżnieniu od rozwiązań opartych na regułach progowych (Geotab, Automatic, eco:Drive)
oraz od prac badawczych, w których klasyfikacja wykonywana jest w środowisku offline,
poza samym pojazdem (Meiring i Myburgh, Bejani i Ghatee, Khan i in.).

Po trzecie, informacja zwrotna przekazywana jest kierowcy w czasie rzeczywistym, w
trakcie samej jazdy, a nie po jej zakończeniu (jak w przypadku eco:Drive) ani w formie
zbiorczego podsumowania tygodniowego (jak w przypadku Automatic). W przeciwieństwie do
rozwiązań korzystających z liczbowego wskaźnika (DriveWell, Geotab) proponowany system
wykorzystuje przy tym metaforę wizualną --- rosnące, zdrowe drzewko przy jeździe
ekonomicznej i więdnące przy jeździe agresywnej --- zgodną z zasadami nudgingu, które
zakładają, że nagradzanie pożądanych zachowań jest skuteczniejszą formą wpływania na
nawyki niż komunikaty o błędach~\cite{thaler2008}, oraz z zasadami gamifikacji, rozumianej
jako wykorzystanie elementów typowych dla gier w kontekście niezwiązanym z rozrywką w
celu zwiększenia zaangażowania użytkownika~\cite{deterding2011}. Po czwarte, cała
funkcjonalność systemu jest dostępna niezależnie od łączności sieciowej ani istnienia
zewnętrznej usługi w czasie --- w odróżnieniu od systemu Automatic, którego całkowite
wyłączenie w 2020 roku pozbawiło funkcjonalności wszystkie sprzedane wcześniej
urządzenia~\cite{automaticlabs}, sprawność proponowanego rozwiązania zależy wyłącznie od
stanu technicznego samego urządzenia zamontowanego w pojeździe.

\subsubsection{Przewagi względem opisanych rozwiązań}

Wynikające z powyższych cech przewagi proponowanego rozwiązania względem opisanych
wcześniej systemów zestawiono w tabeli~\ref{tab:przewagi}. Warto przy tym podkreślić, że
żadna z wymienionych przewag nie jest odosobniona --- pojedyncze elementy proponowanego
podejścia (dane z magistrali CAN, klasyfikacja metodą uczenia maszynowego, gamifikowana
informacja zwrotna, praca bez łączności sieciowej) pojawiają się osobno w
poszczególnych opisanych wcześniej rozwiązaniach lub pracach naukowych, jednak żadne
z nich nie łączy wszystkich tych elementów jednocześnie w ramach jednego, gotowego do
montażu w pojeździe urządzenia.

\begin{table}[!h]
    \centering
    \caption{Przewagi proponowanego rozwiązania względem opisanych alternatyw.}
    \label{tab:przewagi}
    \begin{tabular}{| p{3.4cm} | p{5.8cm} | p{5.3cm} |} \hline
        \textbf{Cecha} & \textbf{Rozwiązania i~badania opisane w~podrozdziałach~\ref{sec:przeglad} i~\ref{sec:uzasadnienie-metod}} & \textbf{Proponowane rozwiązanie} \\ \hline\hline
        Źródło danych & czujniki smartfona (niepewność pomiarowa) lub ograniczony zestaw parametrów OBD & pełny dostęp do sygnałów magistrali CAN pojazdu \\ \hline
        Metoda oceny stylu jazdy & reguły progowe (Geotab, Automatic, eco:Drive) lub model badawczy działający offline (prace naukowe) & model uczenia maszynowego uruchamiany na urządzeniu pokładowym \\ \hline
        Moment przekazania informacji zwrotnej & po zakończeniu jazdy (eco:Drive) lub podsumowanie tygodniowe (Automatic) & w~czasie rzeczywistym, w~trakcie jazdy \\ \hline
        Forma informacji zwrotnej & liczbowy wskaźnik (DriveWell, Geotab) lub aplikacja mobilna wymagająca uwagi kierowcy (EcoGO, Vitiello i~in.) & zgamifikowana wizualizacja na~zintegrowanym z~pojazdem ekranie \\ \hline
        Zależność od łączności/chmury & wymagana stała łączność sieciowa (DriveWell, Geotab, Automatic) & brak zależności --- praca w~pełni lokalna \\ \hline
        Trwałość działania urządzenia & zależna od~istnienia usługi dostawcy (przykład: wycofanie usługi Automatic w~2020~r.) & zależna wyłącznie od~sprawności technicznej urządzenia \\ \hline
        Model kosztowy & jednorazowy koszt sprzętu i/lub cykliczna opłata abonamentowa za~usługę chmurową & jednorazowy koszt sprzętu, bez~opłat cyklicznych \\ \hline
    \end{tabular}
\end{table}

\subsection{Podsumowanie}

Przeprowadzony przegląd istniejących rozwiązań oraz analiza literatury naukowej
pokazały, że rynek systemów analizy stylu jazdy oraz towarzysząca mu literatura
naukowa rozwijają się intensywnie,
jednak w kilku niezależnych kierunkach: rozwiązania komercyjne skupiają się na
dostępności i niskim koszcie wdrożenia kosztem jakości danych i lokalności
przetwarzania, rozwiązania flotowe i ubezpieczeniowe adresowane są do podmiotów
zarządzających, nie do samego kierowcy, prace naukowe nad klasyfikacją metodami
uczenia maszynowego pozostają w większości na etapie badawczym, a systemy stosujące
gamifikację informacji zwrotnej opierają się na danych z czujników smartfona zamiast
na danych pojazdowych. Żadne z omówionych rozwiązań nie łączy jednocześnie wysokiej
jakości danych z magistrali CAN, klasyfikacji metodą uczenia maszynowego oraz
zgamifikowanej informacji zwrotnej dostępnej w czasie rzeczywistym, w pełni lokalnie,
bez zależności od zewnętrznej infrastruktury. Dotyczy to również wymiaru czasowego
informacji zwrotnej --- większość opisanych systemów przekazuje ocenę stylu jazdy
dopiero po zakończeniu przejazdu lub cyklicznie (np.\ raz w tygodniu), co ogranicza jej
praktyczną użyteczność jako narzędzia bieżącej korekty zachowania za kierownicą.

Warto przy tym zestawić rząd wielkości korzyści wynikających z samej poprawy stylu
jazdy, niezależnie od zastosowanej technologii: redukcja zużycia paliwa wynikająca
z ograniczenia agresywnej jazdy sięga, według przywołanych badań, 10--40\% w ruchu
miejskim~\cite{thomas2017}, a rzeczywiste wdrożenia systemów informacji zwrotnej ---
takie jak Fiat eco:Drive czy platforma DriveWell --- potwierdzają, że przekazanie
kierowcy odpowiednio zaprojektowanej informacji zwrotnej pozwala uzyskać wymierną
poprawę konkretnych wskaźników, odpowiednio: średnio 6-procentową redukcję zużycia
paliwa~\cite{fiatecodrive} oraz 20--35-procentową redukcję częstości zdarzeń
niebezpiecznych i rozproszenia uwagi kierowcy~\cite{cmtelematics}. Skala tych efektów
potwierdza, że wybór metody klasyfikacji stylu jazdy oraz formy informacji zwrotnej nie
jest kwestią czysto techniczną, lecz ma bezpośrednie przełożenie na mierzalne korzyści
ekonomiczne i związane z bezpieczeństwem, co uzasadnia szczegółowe podejście do obu
tych zagadnień przyjęte w niniejszej pracy.

Zidentyfikowana w ten sposób luka stanowi uzasadnienie dla podjęcia tematu niniejszej
pracy. W dalszych rozdziałach przedstawiono kolejno: szczegółowe wymagania projektowe
stawiane opracowywanemu systemowi (rozdział~\ref{sec:wymagania}), wprowadzenie do
zasad działania magistrali CAN niezbędne do zrozumienia dalszej części pracy
(podrozdział~\ref{sec:can}), a następnie pełny opis koncepcji, architektury oraz
implementacji proponowanego rozwiązania --- od akwizycji danych z pojazdu, przez
przygotowanie zbioru danych i trening modelu klasyfikującego styl jazdy, po
implementację systemu na docelowej platformie sprzętowej oraz projekt interfejsu
użytkownika (rozdział~\ref{sec:architektura} i kolejne).
